\chapter{Resumen}

En el contexto del rápido aumento del uso de la inteligencia artificial en el diseño de experiencias de usuario, surgen inquietudes importantes sobre la falta de transparencia y la presencia de sesgos en los sistemas algorítmicos que toman decisiones de alto impacto social. Esta memoria aborda esos problemas a partir de un caso real: el Sistema de Admisión Escolar (SAE) de Chile, que asigna estudiantes a establecimientos educacionales mediante un algoritmo de Aceptación Diferida que, si bien es formalmente correcto y está regido por la Ley de Inclusión Escolar N.\textsuperscript{o}~20.845, no comunica su lógica a las familias que lo usan, generando una percepción extendida de arbitrariedad.

El trabajo se apoya en un diagnóstico heurístico de calidad web del SAE real —realizado por el equipo del proyecto Fondecyt N.\textsuperscript{o}~1250492 de la Universidad de Chile, que reveló un cumplimiento de 51\,\% en el sitio informativo y 61\,\% en la plataforma de postulación— y en una revisión sistemática de 96 publicaciones académicas sobre transparencia algorítmica con evaluación de usuarios finales (2017--2025). A partir de ese diagnóstico y de esa revisión se definió un perfil de usuario objetivo (Daniela González) y se construyó, en sucesivas iteraciones documentadas entre mayo y julio de 2026, un prototipo de interfaz que aplica principios de divulgación progresiva, explicabilidad contextualizada y controles interactivos, evolucionando desde un prototipo HTML autocontenido hasta una aplicación web funcional construida con React.

Una verificación interna de trazabilidad frente a la matriz de problemas del diagnóstico registra un cumplimiento del 100\,\% de los puntos aplicables a nivel de código. Una revisión adicional del sitio oficial del SAE, realizada como parte de este trabajo, mostró que solo dos de las seis brechas prioritarias identificadas —búsqueda y encontrabilidad, e inclusión parcial— han sido cerradas de forma independiente por el sistema real, mientras que la transparencia algorítmica interactiva y la interoperabilidad con ClaveÚnica permanecen sin resolver, lo que confirma la vigencia del problema abordado.

\textbf{Esta memoria se encuentra en una etapa intermedia de desarrollo.} El prototipo aún no ha sido sometido a una reevaluación heurística formal ni a una prueba de usabilidad con usuarios reales; ambas validaciones están diseñadas como un plan de cinco fases (Capítulo~3) pero su ejecución, junto con la discusión y las conclusiones que de ella se deriven, queda como trabajo pendiente. Este documento presenta, por tanto, el diagnóstico, el marco teórico, la metodología completa y el estado actual del desarrollo, dejando explícitamente abiertos los capítulos de Resultados de validación, Discusión y Conclusiones hasta que dicha evidencia esté disponible.
